% 第 60 章　对称为何会破缺
\chapter{对称为何会破缺}

\step{反常识开场}

\begin{mainline}
找一支铅笔，把笔尖朝下，小心地竖在光滑的桌面上。手松开之前，先回答一个问题：它会往哪个方向倒？

桌面是平的，重力竖直向下，东南西北没有任何一个方向比别的方向更特别。支配铅笔的规律——牛顿力学和万有引力——对水平面内的所有方向一视同仁。可是铅笔一旦倒下，就必定选中了某一个方向。是谁替它选的？

再看一块冰箱贴。铁原子内部的小磁针像无数枚微型指南针，而电磁学规律同样不偏爱空间中的任何方向。可这块磁铁偏偏坚定地分出南北两极，能吸住便签。规律没有方向，磁铁却有方向。

我们在第 57 章说过，对称就是“做了某种改变之后规律不变”；在第 58 章又看到，对称性会通过诺特定理产生守恒律。既然对称如此强大，这个世界为什么不是一锅完全均匀、毫无结构的粥？星系有旋臂，雪花有六瓣，磁铁有南北，铅笔有倒向——这些“偏爱”究竟从哪里来？

这一章要回答的正是：对称的规律，为什么会生出一个不对称的世界。答案会分三步走出来：先看清“对称”住在哪一层，再看对称状态如何变得待不住，最后看温度、压力和载荷这些参数怎样充当破缺的开关。走完这三步，你会发现铅笔、磁铁、雪花、液晶屏幕，乃至基本粒子的质量，讲的是同一个故事。
\end{mainline}

\step{先猜一猜}

\begin{mainline}
在往下读之前，请把三个预测写下来，越具体越好。写下预测的意义在于：直觉只有被明确说出来，才有可能被实验正面反驳；含糊的“大概吧”永远立于不败之地，也就永远学不到东西。

第一，如果你把同一支铅笔竖立一百次，倒下的方向会均匀散布在各个方向，还是会集中倒向某一边？

第二，设想一支绝对完美的铅笔：笔尖是数学意义上的点，桌面绝对水平，空气完全静止，没有任何振动。它会永远立在那里吗？

第三，把一块磁铁加热到通红，再让它冷却，它的磁性会有什么变化？冷却之后，南北极还会指向原来的方向吗？

大多数人的直觉是：第一题答案“大体均匀”，第二题答案“会永远立着”，第三题答案“磁性消失，冷却后方向随缘”。有意思的是，这三条直觉一条基本正确，一条在经典力学里勉强正确但在真实世界里错误，还有一条藏着一个温度上的“分水岭”。读完本章，你可以回来给自己的直觉打分。
\end{mainline}

\step{动手实验}

\begin{experiment}[竖立铅笔一百次（二十次也够）]
\textbf{材料：}一支铅笔或一根筷子、一张白纸、一支笔。

\textbf{步骤：}
\begin{enumerate}[leftmargin=2em]
  \item 在白纸中央画一个小圆圈，再画出东、南、西、北等八个方向的刻度线，像一张罗盘。
  \item 把铅笔竖在圆圈中央，尽量让它垂直，然后松手，记录它倒下的方向。
  \item 重复至少二十次。每次都尽量用同样的方式竖立，不要故意偏向某一边。
  \item 在罗盘上给每次倒向画一个点，数一数各个方向的点数。
\end{enumerate}

\textbf{你会看到：}单次实验完全无法预测方向；但点的散布大致均匀，没有哪个方向明显占优。偶尔你会遇到“立着不倒”的情况——那通常是因为笔尖变钝或桌面有微小凹痕，铅笔卡进了一个浅浅的小坑里。

\textbf{追加一步：}用两本书把铅笔夹住，只让它能在一条竖直狭缝里倾倒。调节夹持的松紧，体会一个事实：铅笔立得越接近正中央，它待住的时间越长，可一旦开始倒，就倒得越来越快。记住这个感觉，它就是本章的数学图像。

\textbf{追加二步：}拿一把塑料直尺，一端抵在桌面上，另一端用手指竖直向下压。慢慢加力，注意两件事：尺子变弯之前，你能感到它“硬扛”了很久；一旦超过某个力度，它会突然向某一侧弹出，而且你很难预先说出是哪一侧。多试几次，体会“规则不选边，结果必选边”。
\end{experiment}

\begin{mainline}
这个小实验里有三个值得记下的观察。其一，单次结果不可预测；其二，大量结果的统计分布是对称的；其三，“正中央”这个位置并非不存在，而是待不住。第三条正是整章的钥匙。
\end{mainline}

\step{旧模型遇到麻烦}

\begin{mainline}
现在我们手里有几个看似都合理的旧观念，把它们摆在一起，麻烦就来了。

旧观念一：对称的原因产生对称的结果。如果规律对所有方向一视同仁，那么结果也应该没有偏爱。铅笔否定了它——原因对称，结果却不对称。

旧观念二：结果不对称，一定是外界先不对称。也许是风，也许是桌面不平，也许是手抖。这个解释听起来无懈可击，但麻烦在于：无论你把扰动压到多小——抽走空气、用激光校准桌面、用机械手释放——铅笔照样会倒，而且倒向仍然随机均匀。扰动决定的是“哪一次倒向哪边”，却解释不了“为什么非倒不可”。

旧观念三：第 58 章说过，旋转对称意味着角动量守恒。铅笔倒下时分明朝一个方向转了起来，角动量难道凭空出现了？仔细算一笔账：铅笔倒下时确实获得一个很小的角动量，但它与铅笔倾倒方向相关，来源是倾倒前那一点点随机偏移被不断放大；把铅笔和地球看成一个系统，总角动量仍然守恒。而且多次重复实验，倒向均匀分布，平均角动量是零。守恒律没有被违反，被违反的是一个更强的假设——“每一个具体结果都必须和定律一样对称”。

还有一个更隐蔽的旧观念四：只要初始条件给得足够精确，结果就能完全确定，所以“不可预测”只是因为我们测不准。这条在经典力学内部是对的，但它掩盖了一个定量上的绝望：本章后面会估算，铅笔倒下的等待时间只随初始偏移的对数变化——初始偏移每缩小一万倍，等待时间只多出不到一秒。把测量精度提高一亿倍，换来的可预测时间寥寥无几。不稳定系统对初始条件的态度不是“斤斤计较”，而是“无限放大”，这让“精确给定初始条件”在实际中变成一句空话。

真正的麻烦在于，我们一直把两种东西混为一谈：\textbf{定律的对称}和\textbf{状态的对称}。方程旋转不变，推不出方程的每个解都旋转不变。解出一个矛盾，往往需要先发明一个能区分这两个层次的概念。
\end{mainline}

\step{发明新概念}

\begin{mainline}
让我们像物理学家那样，逐个把新概念造出来。

\textbf{第一把刀：定律对称，状态可以不对称。}写出支配系统的方程时，方程本身有旋转对称；但方程允许许多个解，其中每个具体的解完全可以不对称。竖直倒立的铅笔是方程的一个解——它对称，却站不稳；倒向东边的铅笔也是一个解——它不对称，却躺得安稳。世界不住在方程里，世界住在解里。

\textbf{第二个概念：不稳定平衡。}山顶上放一颗球，它可以停在那里——受力确实平衡——但只要偏出头发丝粗细的一点，偏离就会被越放越大。对称的位置如果恰好是山顶，那它在数学上存在，在物理上等于不存在。竖直的铅笔、居里温度以上的磁铁磁化方向，都是这种“存在但待不住”的对称状态。

\textbf{第三个概念：简并。}铅笔倒下之后，倒向东、南、西、北的能量完全一样低。最低能量的状态不是一个，而是一整套彼此等价的状态，物理学家称之为简并。系统必须“选”其中一个躺进去，但它没有偏好的理由——选谁都一样。

\textbf{核心定义：自发对称性破缺。}当一个系统的实际状态——通常是能量最低的状态——不具备支配它的定律所拥有的全部对称性时，我们就说对称性\textbf{自发破缺}了。“自发”二字是关键：没有任何外力替它做决定，不对称是从不稳定性里自己长出来的。

\textbf{对照组：显式破缺。}如果你用手把铅笔推倒，或者用外磁场把一块铁磁化，那是外界先把规矩破了，系统只是服从。这种叫显式破缺，它平淡无奇，本章关心的是没有“推手”的情形。

\textbf{一个度量：序参量。}破缺之后的状态需要一个量来刻画“选了什么、选得多坚决”。对磁铁，这个量是磁化强度——单位体积内小磁针整齐排列的程度。对称相里它为零，破缺相里它非零，方向和大小都有意义。这类量统称序参量。

\textbf{一个裁判：温度与相变。}磁化不是命运的安排，而是温度的判决。把铁加热到约 \(1043\ \mathrm{K}\)（约 \(770\ ^{\circ}\mathrm{C}\)，称为居里温度）以上，热运动的翻腾压过了小磁针彼此对齐的倾向，磁性消失，对称恢复；降温穿过这一点，对称状态重新变得不稳定，磁性自发出现。对称与破缺之间隔着一条明确的温度分界线，这就是相变。

\textbf{最后一对区分：离散破缺与连续破缺。}立一枚硬币，松手，它只有正面朝上、反面朝上两种结局——破缺前后的选择是“离散”的，像岔路口只有两条岔路。立一支铅笔，倒下方向却可以铺满整个圆周——选择是“连续”的。这个差别不只是数量问题：后面会看到，连续对称破缺会附赠一种离散破缺没有的东西——一种几乎不花能量的缓慢波动。磁铁里的自旋波、固体里的声波都是它的化身，硬币却没有对应的“软模式”。

还有两个例子帮你把概念铺开。掰一块圆形苏打饼干，它总沿某一条裂缝断开：饼干的材料和受力大体各向同性，裂缝的方向却是一次具体的选择——只不过材料里预制了许多薄弱点，让选择显得不那么自由。再看《格列佛游记》里的笑谈：小人国为鸡蛋该从大头还是小头打破爆发战争。从哪头打破在“规律”上完全等价，可一旦全社会选了某一头，选择就固化成传统——这是人类社会版本的“状态破缺、定律对称”。

最后看一个生活里的比喻。圆桌宴会上，每个人面前左右各有一杯水，按规矩每个人都只能拿“属于自己”的那一杯，但规则本身分不清左右。只要第一个人伸手拿了左边那杯，全场都会跟着拿左边——方向就这么定了。\textbf{它像什么：}规则完全对称，结果却破缺，而且第一个微小动作决定了全局的方向。\textbf{它不像什么：}宴会上“第一个人”有真实的偏好和手臂，是真正的外因；而物理里的自发破缺不需要任何有偏好的推手，对称状态自己就是不稳定的。比喻到此为止，别把宴会当成物理机制。
\end{mainline}

\step{一句话模型}

\begin{oneline}
规律不偏不倚，但世界只住在“解”里：对称的方程允许一整族彼此等价的不对称解，任何一个真实状态只能挑其中一个住下——于是对称没有消失，而是从“每个状态”退隐到了“所有可能状态的集合”之中。
\end{oneline}

\begin{mainline}
把这句话拆成三截对照前面的例子。方程对称：旋转桌面，铅笔的方程不变。解的集合对称：所有可能的倒向组成的圆周，旋转不变。单个解不对称：任何一次具体的倒下，都指向一个确定方向。自发对称性破缺说的不是“对称没了”，而是“对称换了住所”。

可以用第 57 章的语言再说一遍：物理学关心的对称从来不是“这个东西看起来均不均匀”，而是“做了某种变换之后，规律变不变”。破缺之后做旋转变换，一个解变成另一个解，而所有解的全体和支配它们的规律原封不动。检验对称的探针，从来不指向单个状态，而指向“规律加上全部可能状态”这个更大的对象。明白了这一点，本章的谜面就自动解开：状态的不对称和定律的对称，从一开始就不在同一个层次上，谈不上谁推翻谁。
\end{mainline}

\step{数学翻译器}

\begin{mainline}
现在把“山顶”和“谷底”的感觉画成图。用一个变量 \(x\) 代表系统的状态：对铅笔，它可以表示倒向的偏移；对磁铁，它就是序参量磁化强度。再画一条曲线 \(V(x)\)，表示系统处在状态 \(x\) 时的势能。系统像一颗在曲线上滚动的小球，总是往低处去。

对称相的曲线是一口正中央最低的锅，\(x=0\) 是唯一的谷底，对称状态安稳。破缺相的曲线在中央鼓出一个包，谷底裂成对称的两个：中央仍在原地，却从锅底变成了山顶。一个能写出这两种行为的简单式子是
\[
V(x)=a\,x^{2}+b\,x^{4}\qquad (b>0).
\]
左边回答“状态 \(x\) 的能量有多高”；右边第一项让中央抬高或压低，第二项保证两边远处能量急剧升高，系统不会跑到无穷远。

立刻用特殊情形检查它。若 \(a>0\)，\(x^{2}\) 项占主导，唯一的谷底在 \(x=0\)——对称相。若 \(a<0\)，中央被顶成小山，谷底出现在 \(x=\pm\sqrt{-a/(2b)}\) 处——两个等价的破缺状态。若 \(a=0\)，谷底恰好退回中央，曲线在底部被压得又平又缓，这正是相变发生的临界点。若令 \(b\) 也变成零或负值，\(V(x)\) 会朝 \(|x|\) 增大的一侧无限下跌，系统失去稳定的谷底——公式失效，说明真实系统还需要更高次的项来兜底。而 \(a\) 从正值连续变到负值时，谷底位置 \(\sqrt{-a/(2b)}\) 从零连续长出，破缺是连续发生的——这正对应磁铁在居里温度附近磁化从零连续升高的实验事实。

这幅能量地形图远比看上去通用。再举两个例子，把“参数变号、中央失稳”这一骨架认出来。

\textbf{工程里的破缺：压杆失稳。}拿一把塑料直尺，竖直抵在桌面上，从上方向下压。压力小时，尺子笔直——笔直是对称的解；压力超过某个临界值，尺子突然向某一侧弯出。向左弯和向右弯能量完全相同，方程不偏爱任何一侧，是微小的初始弯曲决定结局。结构工程师把这个现象叫屈曲：脚手架立杆、桥梁压杆、火箭燃料箱的薄壳，都有一条“笔直解失稳”的临界载荷，设计时必须留足余量。这里的参数 \(a\) 由载荷控制：载荷加大，相当于 \(a\) 由正转负，笔直从谷底翻成山顶。

\textbf{屏幕里的破缺：液晶显示。}液晶分子像无数根短小的“分子铅笔”。温度高时它们指向杂乱，各个方向等价；冷却到相变点以下，它们自发地排成大致共同的取向——旋转对称破缺了，序参量就是分子取向。工程师的巧思在于：这个取向可以被电场随时扭转。你手机屏幕上每个像素，本质上就是一个被人为控制的破缺相——电场让分子集体“换方向躺倒”，改变偏振光的通过量，像素随之明暗。每天点亮屏幕，你都在几十万个像素上反复操纵自发对称性破缺。
\end{mainline}

\begin{center}
\begin{tikzpicture}[scale=1.0]
  % 左图：单阱（对称相）
  \draw[->] (-2.2,0) -- (2.2,0) node[right] {\(x\)};
  \draw[->] (0,-0.3) -- (0,2.2) node[above] {\(V\)};
  \draw[thick,domain=-1.55:1.55,samples=50] plot (\x,{0.6*\x*\x});
  \fill (0,0) circle (2pt);
  \node at (0,-0.85) {\small 对称相：谷底在正中央};
  % 右图：双阱（破缺相）
  \begin{scope}[xshift=6.6cm]
  \draw[->] (-2.2,0) -- (2.2,0) node[right] {\(x\)};
  \draw[->] (0,-1.1) -- (0,2.2) node[above] {\(V\)};
  \draw[thick,domain=-1.75:1.75,samples=80] plot (\x,{-1.2*\x*\x+0.35*\x*\x*\x*\x+0.6});
  \fill (-1.31,-0.43) circle (2pt);
  \fill (1.31,-0.43) circle (2pt) node[below right=-1pt] {\small 两个等价谷底};
  \draw[fill=white,thick] (0,0.6) circle (2pt);
  \node at (0.05,1.05) {\small 山顶};
  \node at (0,-1.75) {\small 破缺相：中央失稳};
  \end{scope}
\end{tikzpicture}
\end{center}

\begin{mainline}
\textbf{一次带真实数据的估算：铅笔多久倒下？}竖立的铅笔是一个倒立摆，顶端一旦偏离竖直线一个小角度，偏移就会被指数式放大，放大的特征时间约为 \(\tau\approx\sqrt{l/g}\)。取笔长 \(l\approx 18\ \mathrm{cm}\)，\(g\approx 9.8\ \mathrm{m/s^{2}}\)，得 \(\tau\approx\sqrt{0.18/9.8}\approx 0.14\ \mathrm{s}\)。假设松手时顶端只剩 \(1\ \mathrm{mm}\) 的偏移，要放大到 \(10\ \mathrm{cm}\) 才看得出明显倾倒，需要放大 \(100\) 倍；指数放大 \(100\) 倍约需 \(\ln 100\approx 4.6\) 个特征时间，即约 \(0.6\ \mathrm{s}\)。这和实验手感吻合：竖得越正，等待越久；而一旦偏移进入肉眼可见的范围，剩下的过程快得来不及扶。注意结论的结构：等待时间只随初始偏移的对数增长，所以无论你把铅笔竖得多完美，也只能买来零点几秒的“寿命”。完美主义在这里回报极低。
\end{mainline}

\begin{mathbridge}
为什么谷底恰好在 \(x=\pm\sqrt{-a/(2b)}\)？对 \(V(x)=ax^{2}+bx^{4}\) 求导：
\[
\frac{dV}{dx}=2ax+4bx^{3}=2x\,(a+2bx^{2})=0,
\]
解为 \(x=0\) 或 \(x^{2}=-a/(2b)\)。后者只在 \(a<0\) 时有实数解。再看二阶导数判断稳定性：\(V''(0)=2a\)，所以 \(a>0\) 时中央是谷底，\(a<0\) 时中央翻成山顶——一次求导就把“相变”翻译成了参数 \(a\) 变号。

把 \(x\) 换成平面上的矢量 \((x,y)\)，让势能只依赖到中心的距离：
\[
V=a\,(x^{2}+y^{2})+b\,(x^{2}+y^{2})^{2}.
\]
\(a<0\) 时，谷底不再只是两个点，而是一整个半径为 \(\sqrt{-a/(2b)}\) 的圆环——因为旋转对称是连续的，破缺后的等价基态也排成连续的一圈。这幅图形的剖面像一顶墨西哥草帽。沿圆环方向移动不花任何能量，只有沿半径方向爬出谷底才花能量。这个观察通向一个重要结论（戈德斯通定理的直觉版）：每当一个连续对称自发破缺，就必然存在一种“沿谷底慢慢晃”的、几乎不花能量的波动模式。磁铁里的自旋波、固体里的声波，都是这类模式的化身。
\end{mathbridge}

\begin{center}
\begin{tikzpicture}[scale=1.1]
  % 墨西哥帽剖面 + 谷底圆环俯视
  \draw[thick] plot[smooth] coordinates {(-2,1.55) (-1.2,0.15) (-0.45,0.88) (0,1.02) (0.45,0.88) (1.2,0.15) (2,1.55)};
  \draw[dashed] (0,0.15) ellipse (1.2 and 0.3);
  \draw[->] (0,1.02) -- (0,2.3) node[above] {\(V\)};
  \fill (-1.2,0.15) circle (1.6pt);
  \fill (1.2,0.15) circle (1.6pt);
  \draw[fill=white,thick] (0,1.02) circle (1.6pt);
  \node at (1.7,0.62) {\small 谷底连成一圈};
  \draw[->] (1.8,0.5) -- (1.38,0.24);
  \node at (0,-0.7) {\small “墨西哥帽”：沿虚线圆环走不花能量};
\end{tikzpicture}
\end{center}

\step{模型的边界}

\begin{mainline}
一个好模型必须交代自己在哪里会失效。本章模型有六条边界，前三条回答“它什么时候说得对”，后三条提醒“它容易被怎样讲错”。

\textbf{边界一：破缺的是状态，不是定律。}铅笔倒下后，牛顿定律照样旋转对称；磁铁磁化后，麦克斯韦方程照样不偏爱方向。守恒律一条也没有被破坏。如果有人在破缺的世界里宣称“对称性被推翻了”，那是把地图撕了却以为领土变了。

\textbf{边界二：一个反例，专治“只要够完美就不会破缺”的直觉。}你可能会想：竖立铅笔会倒，无非是因为做不到绝对完美；若真有一支理想铅笔、理想桌面、真空环境，它就永远立着，所谓自发破缺不过是外界扰动的遮羞布。这个反驳在经典力学里有一半道理——理想倒立摆确实是一个严格的解。但它在真实世界站不住脚，原因不是手艺不够，而是量子力学的基本原理：一个状态不可能同时具有严格确定的位置（恰好竖直）和严格确定的动量（恰好静止），这是不确定性关系对“状态”本身的限制，不是测量仪器粗糙。所以“毫无扰动的完美竖立”这个状态在物理上根本不存在，破缺不需要外界的借口。

\textbf{边界三：系统不够大时，对称会回来。}考虑一块只有几百个原子的小磁体。它虽然选了一个磁化方向，但等上足够长的时间，热涨落或量子隧穿能让整个磁化方向翻转过去再翻转回来；平均而言，它对任何方向都没有偏爱。严格的数学事实是：有限大小的系统，真正的基态常常是所有方向的对称叠加，自发破缺只在“系统趋于无限大”的极限下完全成立。宏观磁铁包含约 \(10^{23}\) 量级的原子，翻转等待时间远远超过宇宙年龄，所以我们放心地说它的对称破缺了。“破缺”是一个关于极限的概念，日常语言里的“永远”在这里是“等到宇宙热寂也等不到”的缩写。

\textbf{边界四：不要把希格斯机制讲过头。}第 59 章讲了规范对称性，那里的故事在本章迎来续集：现代理论认为，充满整个空间的希格斯场发生自发破缺，从对称状态“倒进”某个真空值，于是 \(W\)、\(Z\) 玻色子和夸克、电子等基本粒子获得质量项，光子则保持无质量。这是希格斯机制的概念角色——它解释了基本粒子质量参数的由来，解释了电与弱相互作用为何表现迥异。但它不是“万物的称重机”：你身体质量的绝大部分来自质子和中子，而质子、中子质量的约九成九来自夸克和胶子的相互作用能，与希格斯场基本无关。说“希格斯场赋予了所有物质质量”，就把一个精确的概念角色吹成了一个漂亮的误会。

\textbf{边界五：不要把自发破缺当成解释一切不对称的万金油。}你的心脏长在左边、大多数人习惯用右手、地球上的一条河总是侵蚀某一侧河岸——这些不对称各有各的具体成因，多数是历史偶然被固化，与“对称状态失稳”无关。自发破缺是一个有严格条件的概念：定律（近似）对称、对称状态失稳、存在一族等价状态、系统足够大。四条缺一条，就别急着套用它。还要警惕反向误用：看见世界的不对称就反推“规律一定有偏好”，这正是本章开头拆掉的旧观念。

\textbf{边界六：破缺的方向可以约定，破缺的存在不能。}磁铁的南北极叫什么名字、电池哪头标正号，是人类的约定；但“磁铁必然选出某一个方向”不是约定，是稳定性决定的。混淆这两层，就会把物理事实错当成文化习惯，或把文化习惯错当成物理定律。
\end{mainline}

\begin{challenge}
\textbf{有限系统的对称恢复与“薛定谔猫”的联系。}在双阱势里放一个量子粒子，它的最低能量状态不是“住在左边”或“住在右边”，而是左右两个位置的叠加——隧穿效应把两个谷底连通了，真正的基态恢复了对称。只有当系统自由度趋于无穷，左右两个状态之间的隧穿概率被指数压低到零，叠加才实际等同于“选定一边”。宏观磁铁的磁化方向之所以稳定，正是因为它是天文数字个微观自由度“集体站队”的结果。这也解释了为什么“破缺”总是和“宏观”“凝聚”“大量自由度”同时出现。

\textbf{戈德斯通定理（陈述版）。}若一个连续对称性自发破缺，而相互作用又是短程的，则必然出现能量可以任意低的激发模式，称为戈德斯通模式。固体中的声波（声子）、铁磁体中的自旋波都是实例。希格斯机制的特殊之处在于：破缺的是规范对称性，原本该出现的戈德斯通模式被规范场“吸收”，转而成为有质量的规范玻色子的纵向分量——这正是 \(W\)、\(Z\) 变重的机制。本书只陈述角色，不作推导。
\end{challenge}

\step{回头解释开场现象}

\begin{mainline}
回到开场的三个悬念，用新模型重新讲一遍。

\textbf{铅笔。}桌面各方向等价——定律对称；竖直姿态是方程的解，却是不稳定的山顶；松手后，初始那一点点无法消除的偏移被指数放大，约零点几秒后铅笔躺进某个方向的谷底。单次倒向不可预测，因为决定它的是放大前的微小偏移；多次重复方向均匀，因为偏移本身没有偏爱。角动量没有凭空出现：铅笔获得的角动量与它的倒向绑定、方向随机，统计平均为零，铅笔与地球的总角动量守恒。诺特定理不但没有被推翻，反而完成了最漂亮的自我辩护：正因为定律旋转对称，规律才无法“规定”铅笔倒向哪里——方向不是被规律选的，是被涨落选的，规律只负责把涨落放大成定局。

\textbf{磁铁。}高温下，热运动让每个小磁针随机乱指，序参量为零，系统处于对称相；冷却穿过居里温度（铁约 \(770\ ^{\circ}\mathrm{C}\)，蜡烛火焰足以越过它），对称状态失稳，小磁针彼此对齐的倾向被放大，磁化强度自发长成，方向由冷却时刻的涨落或残余磁场决定。再次加热，对称恢复——这就是为什么磁铁烧红再冷却后，南北极常常换了指向。

\textbf{结冰与雪花。}一杯水在液态时对任何方向一视同仁：把水盆旋转任意角度，水看不出区别。冻结之后，水分子排成晶格，选定了几组特定的晶轴方向——连续的旋转对称破缺成离散的六重对称。注意破缺是“降级”而不是“清零”：冰仍然有对称性，但只剩六十度整数倍的旋转。雪花的六瓣花样正是这场破缺的遗迹：微观晶轴的六重对称被生长过程放大成肉眼可见的六重图案。同一场降雪里没有两片雪花相同，因为每片雪花在飘落的旅程中经历的温度湿度略有不同，细枝末节被放大成枝杈差异；但它们共享同一个六重骨架，因为那是定律允许的选择。这个例子把本章三个主角——对称、破缺、放大——串在了一片雪花里。

\textbf{世界的结构。}现在可以回答开场的大问题了：均匀与对称属于“定律”和“高温”，结构与偏爱属于“解”和“低温”。星系、雪花、磁铁、乃至基本粒子的质量谱，都是规律冷却到各种破缺相之后长出来的结构。对称性并没有和丰富性为敌——恰恰相反，正因为有一整族等价的解可供选择，世界才有了“做出选择”的机会。对称规定了什么允许发生，破缺决定了什么真的发生。

这也为下一章埋好伏笔：第 61 章将讨论“基本”与“涌现”两个解释方向。破缺相里的规律——比如固体中的声速、磁畴的边界——不必从微观方程逐字读出，它们是大量自由度集体选择的结果；但微观对称性决定了哪些集体选择被允许。两个方向在“对称与破缺”这里交汇。
\end{mainline}

\step{脑洞实验与挑战题}

\begin{thought}[宇宙是一支冷却的铅笔]
今天的粒子物理图景认为，宇宙极早期温度极高，电磁相互作用与弱相互作用本是一家，由统一的对称性支配。随着宇宙膨胀冷却，充满空间的希格斯场从对称状态“倒进”某个真空值——就像铅笔倒下、磁铁降温——统一的对称自发破缺，\(W\)、\(Z\) 玻色子变重，光子保持无质量，电与弱从此分道而行。

那么问题来了：如果宇宙比我们能观测的范围大得多，不同区域“倒”的方向会不会不一样？理论给出的答案是“可能”。就像一块磁铁里存在指向不同的磁畴、一盆水结冰时不同晶区的晶向彼此错开，早期宇宙的破缺也可能在空间上“选得不一致”，交界处留下被称为畴壁、宇宙弦的“缺陷”结构。寻找这些远古疤痕的踪迹，是现代宇宙学的窗口之一。

这类想法还附带一条深刻的物理教训：破缺一旦发生，就抹去了它自己的部分历史。一块磁畴内部的人无法从本地测量推出“当初若不破缺会怎样”，正如我们无法从今天的磁铁推出它冷却前那个无方向的对称世界。宇宙学观测之所以珍贵，正因为大爆炸的余温里还可能冻着破缺发生瞬间的信息——那是世界“做出选择”时留下的指纹。

\textbf{它像什么：}冷却、失稳、随机选择、留下畴界——和磁铁降温完全同构。\textbf{它不像什么：}宇宙没有“桌面”也没有“手”，破缺的场充满整个空间，不存在任何外部参考物；而且“倒向”是场的内部取向，不是空间里的某个方位。别把宇宙真想成一支立在桌上的铅笔。
\end{thought}

\begin{mainline}
下面四道题重在推理，不靠计算。每题后面附有思路提示，但请先自己想透再看。它们分别对应本章的四个台阶：离散与连续、不稳定性的定量后果、破缺与结构、序参量的实验判据。

\begin{puzzles}
  \item 把一枚硬币立在桌面上松手，它倒下；把一支铅笔立在桌面上松手，它也倒下。两次“对称破缺”有什么本质差别？提示：数一数两个系统各自有多少个等价的结局；回忆连续对称破缺特有的那种“几乎不花能量的缓慢波动”，想想它在硬币情形中是否存在。
  \item 有人说：“一支绝对完美的铅笔，在绝对完美的桌面上，只要没有任何扰动，就永远不会倒，所以自发破缺这个概念是多余的。”请用本章的内容反驳他，要求给出两条独立的理由。提示：一条从“平衡是否稳定”入手，估算等待时间如何依赖初始偏移；另一条从量子力学对“状态”的限制入手，想想“恰好竖直且恰好静止”是不是一个存在的状态。
  \item 液态水对任何方向一视同仁，冰却选定了晶轴方向——对称“变少”了。可为什么偏偏是对称更少的冰，才能长出雪花那样规则而丰富的花样？提示：对称相里系统有几个可供选择的状态？破缺相里呢？“结构”需要的是什么？
  \item 把一根铁钉在火上烧到通红，迅速移近一个指南针（假设不接触、不烫坏指南针），指南针会明显偏转吗？冷却之后再试一次呢？提示：想想序参量在居里温度上下的取值；再想想“没有任何宏观磁性”和“对磁场毫无响应”是不是同一回事——顺磁响应始终存在，只是太弱。
\end{puzzles}
\end{mainline}
